% Resumo e palavras-chave compartilhados pela edicao em portugues.
\newcommand{\capmabstract}{%
A pesquisa e a regulação de cibersegurança automotiva concentram-se no produto:
a ISO/SAE~21434 e a UN~R155 governam o carro, suas ECUs e seu \emph{back end}.
A fábrica que constrói o carro não é governada por nada equivalente, embora
seja o ativo cuja falha interrompe a produção, deixa uma cadeia de fornecimento
continental sem peças e, no evento cibernético mais caro já registrado no Reino
Unido, custou a uma economia cerca de \pounds1,9 bilhão. Este artigo trata a
planta automotiva, e não o veículo, como infraestrutura crítica, e faz uma
pergunta que a literatura centrada no veículo não responde: por quais caminhos
um atacante efetivamente alcança a linha de montagem? Contribuímos com (i)~um
corpus curado de 20 incidentes publicamente reportados na manufatura automotiva
entre 2017 e 2025, cada um reconstruído como um caminho em nível de ativo com
rótulo explícito de confiança; (ii)~o \capm{}, um modelo de caminhos de ataque
que atravessa TI corporativa, o plano de identidade, nuvem e SaaS, a DMZ
industrial, as camadas MES/SCADA/CLP e a interface com fornecedores,
instanciado como uma arquitetura de referência de 12 zonas IEC 62443, 48 ativos
e 91 passos de ataque rotulados com 43 técnicas MITRE \attck{} Enterprise e
\attck{} for ICS; (iii)~uma análise probabilística que combina o ranqueamento
dos $k$ caminhos mais prováveis com uma simulação Monte Carlo de campanhas que
quantifica a perda anual esperada; e (iv)~um catálogo de controles mapeado para
a IEC 62443-3-3 e o NIST CSF 2.0, avaliado por valor marginal, necessidade e
custo-efetividade. Quatro resultados se seguem. Primeiro, o caminho de menor
resistência até a parada de produção nunca toca um controlador: a rota mais
provável é $11{,}6\times$ mais provável do que a melhor rota que alcança o
nível~2 de Purdue ou abaixo, e o corpus concorda: 55\% dos incidentes causaram impacto de disponibilidade,
enquanto apenas 10\% teriam interagido com equipamentos abaixo do nível~3. Segundo, a capacidade de recuperação domina a
prevenção em termos de perda esperada: \emph{backups} imutáveis e testados e um
modo degradado de produção ensaiado respondem pelas duas maiores reduções
obtidas por um único controle, ao passo que a governança de tokens SaaS e a
garantia de segurança de fornecedores são os controles que um programa maduro
menos pode remover. Terceiro, o endurecimento parcial desloca o risco em vez de
removê-lo: um programa voltado a cadeia de fornecimento e nuvem reduz a perda
por espionagem em 73\% e \emph{aumenta} a perda por ransomware em 26\%, ao
empurrar um ator racional para o objetivo de disponibilidade. Quarto, medido contra quatro priorizações de referência sobre o mesmo corpus,
o modelo coloca o incidente observado mediano no percentil $4{,}8$ das rotas até sua consequência,
onde heurísticas de profundidade em OT e de travessia de zonas o colocam nos
percentis $95{,}9$ e $98{,}9$; já uma ordenação que ignora completamente a
detecção tem desempenho indistinguível, de modo que o termo de contenção
justifica seu lugar na estimativa de perda, e não no ranqueamento. O modelo, o
corpus, os experimentos e as figuras são publicados como artefato aberto e sem
dependências externas.%
}

\newcommand{\capmkeywords}{%
infraestrutura crítica, manufatura automotiva, grafos de ataque, tecnologia
operacional, MITRE \attck{} for ICS, IEC 62443, segurança da cadeia de
fornecimento, quantificação de risco%
}
